\documentclass{report}
\usepackage{amsmath,amssymb}
\setlength{\parindent}{0mm}
\setlength{\parskip}{1em}
\begin{document}
\begin{center}
\rule{6in}{1pt} \
{\large Bobby Philip \\
{\bf The Advanced Multi-Physics (AMP) Package}}

OAK RIDGE NATIONAL LABORATORY \\ PO BOX 2008 MS6164 \\ OAK RIDGE TN 37831-6164
\\
{\tt philipb@ornl.gov}\\
Kevin Clarno\\
Rahul Sampath\\
Mark Berrill\\
Srikanth Allu\\
Gary Dilts\end{center}


Simulations of multiphysics systems are becoming increasingly important
in many science application areas.
In this talk we describe the design and capabilities of the Advanced
Multi-Physics (AMP) package designed with multi-domain
multi-physics applications in mind. AMP currently builds powerful and
flexible deterministic multiphysics simulation
algorithms from lightweight operator, solver, linear algebra, material
database, discretization, and meshing
interfaces. AMP�s agnostic framework provides at least two important benefits.

Firstly, it enables domain scientists with considerable investment in physics codes to
experiment with multi-physics couplings without having to adopt dramatically
different data structures for multiphysics problems. Physics components
are represented as discrete operators that map vectors between domain
and range spaces and encode the action of physics operators on vectors. A
minimal interface is required
of each operator, consisting of the ability to specify: the action of the
operator on a vector, operator
(re)initialization, and approximate linearization (optional). Domain
scientists can thus preserve their
investments in existing physics by encapsulating them as operators while
at the same time preserving the
ability to create new physics operators with potentially different data
structures. In future releases, operator
adaptor interfaces will further minimize code modifications required. An
operator can now represent
a single physics or potentially multiple physics either as a composition
of single physics operators or
as a single tightly coupled multiphysics operator. The minimal operator
interface allows easy import of
higher fidelity models without significant code rewrite. This approach
has also naturally allowed for multi-domain
applications. Finally, it allows domain scientists and mathematicians to
study and analyze the strength of
the couplings between various physics components, numerical error
propagation, and the sensitivities of
output quantities of interest to input and model parameters, data,
boundary and initial conditions.

Secondly, considerable research investment by DOE, through the SciDAC
TOPS, Advanced Simulation and
Computing (ASC) projects, and others has led to the development of highly
sophisticated software
for solving individual physics efficiently at the petascale (e.g. PETSc,
Trilinos, SUNDIALS
, Hypre). Moreover, further investments in software for solution methods
for nonlinear and
linear multiphysics problems will have to be made as applications become
more complex and computer
architectures evolve. By carefully specified solution interfaces, AMP is
able to leverage these existing
investments, while keeping the door open to future solution
methodologies. Within AMP, physics solvers
are considered as inverse or approximate inverse operators with a minimal
interface. The concept of
operator composition now again enables AMP users to harness existing
software to efficiently realize
different multiphysics solution algorithms and experiment with new
physics solvers without extensive
code rewrites. Through the AMP solver interfaces a subset of the PETSc,
Trilinos, and SUNDIALS solvers and time
integrators are already available to users with the ability to add new
solvers as required. This allows
users to balance needs for accuracy, robustness, and efficiency while
respecting the nature of couplings
between different physics components, and choosing appropriate solvers
for the individual physics. Using
Jacobian Free Newton Krylov (JFNK) and nonlinear Krylov methods AMP
provides the ability to solve
the multiphysics systems composed from individual physics in a
nonlinearly consistent manner, with the
flexibility to choose the ordering of the individual physics in the
solution process, as well as consider a
variety of alternative multiplicative (i.e., the updated solution in a
module are immediately passed on to the
next coupled module for its solution processing), additive (i.e.,
components of weakly coupled subsystems
are solution processed asynchronously), and hybrid Schwarz solution
algorithms. The same concepts for
operator and solver composition and decomposition are used for
time-dependent problems, which will be
required for the solution of multiphysics modules with time-scales that
differ by orders of magnitudes.
AMP users have the flexibility to choose different time integrators for
the different physics modules,
tailored according to the optimal strategy of each particular model using
composition of individual time
integrators.

AMP was developed with funding from the DOE Office of Nuclear Energy
under the Nuclear Energy Advanced
Modeling and Simulation (NEAMS) Program. In this and related talks we
will describe how AMP has been
used to solve coupled multi-domain multi-physics problems related to energy applications.


\end{document}
